\section*{Capítulo \ref{ch:g12:mountain-building} --- La formación de las montañas}

\begin{solution}{exo:g12:mountain-building:1}
(a) \omterm{def:g12:mountain-building:thrust}{cabalgamiento}; (b) \omterm{def:g12:mountain-building:ophiolite}{ofiolita}; (c) \omterm{def:g12:mountain-building:fold}{sinclinal}; (d) \omterm{def:g12:mountain-building:suture}{sutura}; (e) \omterm{def:g12:mountain-building:fold}{anticlinal};
(f) \omterm{def:g12:mountain-building:nappe}{manto de corrimiento}.
\end{solution}

\begin{solution}{exo:g12:mountain-building:2}
(1) Los flancos buzan alejándose del \omterm{def:g12:mountain-building:fold}{eje del pliegue}: un \omterm{def:g12:mountain-building:fold}{anticlinal}. (2) Las
capas más antiguas. (3) No: la erosión sigue las capas duras y blandas, no la
forma del \omterm{def:g12:mountain-building:fold}{pliegue}; el núcleo de un \omterm{def:g12:mountain-building:fold}{anticlinal} puede quedar excavado en un valle.
\end{solution}

\begin{solution}{exo:g12:mountain-building:3}
(1) $e = 1 - \cos\ang{30} = 1 - 0.866 = 0.134$, alrededor del \qty{13}{\percent}.
(2) $L_0 = 6.0/\cos\ang{30} \approx \qty{6.9}{km}$.
\end{solution}

\begin{solution}{exo:g12:mountain-building:4}
(1) Una \omterm{def:g11:ocean-ridges:normal-fault}{falla normal}; la corteza se estiró. (2) Sí: es una \omterm{def:g12:mountain-building:thrust}{falla inversa} que
buza menos de \ang{45}.
\end{solution}

\begin{solution}{exo:g12:mountain-building:5}
(1) Sedimentos marinos, basaltos almohadillados, complejo de \omterm{def:g11:igneous-rocks:intrusion}{diques} laminares,
\omterm{def:g11:ocean-ridges:gabbro}{gabro}, \omterm{def:g11:ocean-ridges:peridotite}{peridotita}. (2) La lava solo forma almohadillas cuando se enfría en el
agua. (3) Una \omterm{def:g12:mountain-building:ophiolite}{ofiolita} es un fragmento del océano que había entre dos
continentes; cuando chocaron, láminas de su fondo fueron empujadas sobre sus
bordes en lugar de hundirse, así que se encuentran donde se cerró el océano,
a lo largo de la \omterm{def:g12:mountain-building:suture}{sutura}.
\end{solution}

\begin{solution}{exo:g12:mountain-building:6}
(1) $L_0 = 2.1 + 1.8 + 2.4 + 1.6 = \qty{7.9}{km}$;
$e = (7.9 - 6.0)/7.9 \approx 0.24$, el \qty{24}{\percent}. (2) $\cos\theta =
6.0/7.9 \approx 0.759$, $\theta \approx \ang{41}$.
\end{solution}

\begin{solution}{exo:g12:mountain-building:7}
(1) $45\cos\ang{10} \approx \qty{44.3}{km}$. (2) $45\sin\ang{10} \approx
\qty{7.8}{km}$. (3) Esos \omterm{def:g12:mountain-building:fold}{pliegues} acortan en $1 - \cos\ang{12} \approx 0.0219$:
el cinturón tendría que haber medido $44.3/0.0219 \approx \qty{2000}{km}$ de
ancho antes del plegamiento.
\end{solution}

\begin{solution}{exo:g12:mountain-building:8}
(1) La superposición: leída sin la falla, la roca superior se tomaría por la
más joven. (2) $500 - 75 = 425$ millones de años. (3) $\qty{40000}{m} /
\qty{800}{m} = 50$ veces.
\end{solution}

\begin{solution}{exo:g12:mountain-building:9}
(1) $(30.23 - 19.38) \times 111.2 = 10.85 \times 111.2 \approx
\qty{1207}{km}$. (2) $1207/14 \approx \qty{86}{km}$ por millón de años,
\qty{86}{mm}, o unos \qty{8.6}{cm}, al año. (3) \qty{9}{m} por siglo son
\qty{9}{cm} al año: los dos valores coinciden con un margen de un
\qty{5}{\percent}.
\end{solution}

\begin{solution}{exo:g12:mountain-building:10}
(1) $H = 35/(1 - 0.40) \approx \qty{58}{km}$. (2) $\Delta H \approx
\qty{23.3}{km}$; se eleva $23.3 \times 550/3380 \approx
\qty{3.8}{km}$ más, y su \omterm{def:g12:continental-crust:root}{raíz} añadida tiene $23.3 - 3.8 \approx \qty{19.5}{km}$ de profundidad.
\end{solution}

\begin{solution}{exo:g12:mountain-building:11}
(1) El dibujo multiplica $\tan\theta$ por 3: $\tan\theta = \tan\ang{40}/3
\approx 0.280$, $\theta \approx \ang{15.6}$. (2) Real: $1 - \cos\ang{15.6}
\approx 0.037$, el \qty{3.7}{\percent}; leído en el dibujo: $1 - \cos\ang{40}
\approx 0.234$, el \qty{23}{\percent}. (3) Una escala vertical exagerada
aumenta los buzamientos y alarga las capas, así que hace que el \omterm{def:g12:mountain-building:shortening}{acortamiento}
parezca varias veces demasiado grande.
\end{solution}

\begin{solution}{exo:g12:mountain-building:12}
(1) $L_0 = 10/(1 - 0.317) \approx \qty{14.6}{km}$. (2) $1 - \cos\ang{57.5}
\approx 1 - 0.537 = 0.463$, el \qty{46}{\percent}. (3) En el \omterm{def:g12:mountain-building:fold}{pliegue}
redondeado, solo una parte corta de la capa buza \ang{57.5}; cerca de las
charnelas es casi horizontal, y las partes horizontales no se acortan. En el
\omterm{def:g12:mountain-building:fold}{pliegue} de flancos rectos, toda la capa buza \ang{57.5}.
\end{solution}

\begin{solution}{exo:g12:mountain-building:13}
(1) $\sqrt{16.6^2 + 41.8^2} \approx \qty{45.0}{mm/yr}$. (2) $41.8 - 23.5 =
\qty{18.3}{mm}$ al año, es decir, unos \qty{18}{km} de \omterm{def:g12:mountain-building:shortening}{acortamiento}
norte--sur por millón de años al norte de Lhasa. (3) La propia Lhasa se ha
desplazado hacia el norte, y los terremotos del cinturón se extienden por Asia
mucho más al norte del Himalaya: la corteza sigue rompiéndose allí.
\end{solution}

\begin{solution}{exo:g12:mountain-building:14}
(1) Su buzamiento es casi nulo, así que el \omterm{def:g12:mountain-building:shortening}{acortamiento} es el desplazamiento,
\qty{45}{km}. (2) $45/(1 - \cos\ang{12}) = 45/0.0219 \approx \qty{2060}{km}$.
(3) Las grandes convergencias las absorben sobre todo los \omterm{def:g12:mountain-building:thrust}{cabalgamientos}, los
\omterm{def:g12:mountain-building:nappe}{mantos de corrimiento} y los \omterm{def:g12:mountain-building:fold}{pliegues} apretados; los \omterm{def:g12:mountain-building:fold}{pliegues} suaves absorben
muy poco.
\end{solution}

\begin{solution}{exo:g12:mountain-building:15}
(1) $\qty{8848.86}{m} = \qty{8848860}{mm}$ en $\num{5e7}$ años: unos
\qty{0.18}{mm} al año. (2) $40/0.18 \approx 220$: la elevación es más de
doscientas veces más lenta que el movimiento de la India. (3) La caliza se
depositó por debajo del nivel del mar, y la erosión ha eliminado las rocas que
tenía encima: se ha elevado más que su altura actual.
\end{solution}

\begin{solution}{pb:g12:mountain-building:1}
\textbf{1.} $2\pi \times 6371/360 \approx \qty{111.2}{km}$.

\textbf{2.} Entre hace 50 millones de años (\ang{0.68}~S) y hace 45 (\ang{3.30}~N):
hace $50 - 5 \times 0.68/3.98 \approx 49$ millones de años.

\textbf{3.} $(19.38 - 0.68) \times 111.2 = 18.70 \times 111.2 \approx
\qty{2079}{km}$.

\textbf{4.} $2079/16 \approx \qty{130}{km}$ por millón de años: \qty{130}{mm}
al año, unos \qty{13}{cm}.

\textbf{5.} De 66 a 60: $7.28 \times 111.2/6 \approx 135$; de 60 a 55: $4.86
\times 111.2/5 \approx 108$; de 55 a 50: $6.56 \times 111.2/5 \approx 146$; de 50
a 45: $3.98 \times 111.2/5 \approx 89$; de 45 a 40: $2.29 \times 111.2/5
\approx 51$ (\unit{mm} al año). Lo más rápido, entre hace 55 y 50 millones de años.

\textbf{6.} $(20.00 - 5.59) \times 111.2/40 = 14.41 \times 111.2/40 \approx
\qty{40}{mm}$ al año.

\textbf{7.} $40/130 \approx 0.31$: menos de un tercio.

\textbf{8.} Del paso 55--50 (146) al paso 50--45 (89): hacia hace 50
millones de años.

\textbf{9.} La \omterm{def:g12:continental-crust:continental-crust}{corteza continental} es demasiado ligera para hundirse: cuando
llegó a la fosa, la \omterm{def:g11:subduction-zones:subduction}{subducción} ya no podía arrastrar libremente a la India, y
el otro continente la frenó.

\textbf{10.} $(20.00 + 0.68) \times 111.2 = 20.68 \times 111.2 \approx
\qty{2300}{km}$.

\textbf{11.} $(29.65 - 21.15) \times 111.2 = 8.50 \times 111.2 \approx
\qty{945}{km}$.

\textbf{12.} $2300 - 945 \approx \qty{1354}{km}$.

\textbf{13.} Hoy, $(29.65 - 20.00) \times 111.2 \approx \qty{1073}{km}$;
hace 50 millones de años, $(21.15 + 0.68) \times 111.2 = 21.83 \times 111.2
\approx \qty{2427}{km}$ (y $2427 - 1073 = 1354$).

\textbf{14.} $e = 1354/2427 \approx 0.56$, el \qty{56}{\percent}.

\textbf{15.} \omterm{def:g12:mountain-building:fold}{Pliegues} y \omterm{def:g12:mountain-building:thrust}{fallas inversas}: \omterm{def:g12:mountain-building:thrust}{cabalgamientos} que transportan \omterm{def:g12:mountain-building:nappe}{mantos de corrimiento}.

\textbf{16.} $H = 35 \times 2427/1073 \approx \qty{79}{km}$.

\textbf{17.} $\Delta H = 79 - 35 = \qty{44}{km}$; se elevaría
$44 \times 550/3380 \approx \qty{7.2}{km}$ más.

\textbf{18.} $\Delta H = 4.6 \times 3380/550 \approx \qty{28}{km}$, una
corteza de $35 + 28 = \qty{63}{km}$ de espesor, cerca de los \qty{70}{km}
medidos.

\textbf{19.} La corteza no conserva su sección: partes de Asia son expulsadas
hacia el este, y la erosión elimina roca de la parte alta de las montañas.

\textbf{20.} $2079 + 2300 \approx \qty{4380}{km}$: desde el final del
Cretácico, la India ha viajado unos \textbf{\qty{4400}{km} hacia el norte}.
\end{solution}
